% ---
% id: meat-core/rem:cgate
% kind: rem
% title: "Degeneracy of the licensed-class half of the gate"
% summary: >-
%   In either format the licensed-class half of the constraint gate is a
%   genuine type check in trusted-source mode at d_A ≥ 2, vacuous in
%   untrusted-source mode, and at d_A = 1 rejects only what a producer can
%   trivially avoid. In the last two cases the key length rests on the
%   attack-set licence alone.
% prominence: standard
% status: accepted
% uses: [def:homog, prop:trivial-alice-system, def:declaration, thm:keylength, rem:homog, def:feasible, thm:convexity, thm:removal, thm:choi, thm:gradient, def:admissible-declaration]
% mentions: [req:constrained-key-gate, req:degenerate-gate-report, protocol-compiler/rem:party-local-soundness]
% assumes: [A-ATK, A-SRC]
% refs: []
% provenance: newly-derived
% lean:
% ---

The numerical machinery --- Definition~\ref{def:feasible},
Theorems~\ref{thm:convexity}, \ref{thm:removal}, \ref{thm:choi} and
\ref{thm:gradient} --- is valid for arbitrary $\CList$, with
$\hlow(\CList)\ge\hlow(\varnothing)$; what the two-part gate on $\CList$
(Requirement~\ref{req:constrained-key-gate}) decides is only whether its
output may reach the key length of Theorem~\ref{thm:keylength}.

\begin{remark}\label{rem:cgate}
Part (i) of the two-part gate --- ``the list lies in the licensed class
of the declared source mode'' --- is the half a program can check. How
much it checks depends on the source mode and on $d_A$.

\emph{\tref[trusted-source mode]{Trusted-source mode} with $d_A\ge2$:} the restriction is genuine
--- the multiples of $\tau_A$ are a one-dimensional subspace of the
$d_A^2$-dimensional real space of Hermitian operators on $A$, so a
\tref{general} entry is in general not a licensed one --- and a declaration
can fail part~(i). \emph{The check is nevertheless a \emph{type} check}: an
entry declares either its scalar $c_k$ or an operator $\Theta_k$
(Definition~\ref{def:homog}), and part~(i) inspects which, rather than
comparing the operator against $c_k\tau_A$. What is being rejected is a
declaration that \emph{says} it is general, never one whose numbers came out
near a multiple; the class is the same and the decision is exact.

\emph{$d_A=1$:} it is a type check that rejects nothing a producer cannot
trivially avoid. An entry declared \tref{general} fails it, as at every
$d_A$ (Definition~\ref{def:homog}); but every $1\times1$ Hermitian
$\Theta_k$ equals $\Theta_k\cdot[1]=c_k\tau_A$
(Proposition~\ref{prop:trivial-alice-system}\,(iii)), so such an entry could
always have been declared homogeneous with the same condition, and no
declaration \emph{need} fail part~(i).

\emph{\tref[untrusted-source mode]{Untrusted-source mode}, any $d_A$:} it is
\emph{vacuously true}. The entry \emph{type} is scalar by
Definition~\ref{def:declaration}, so part~(i) degenerates from a
mathematical condition to a type check that a declaration passes by
construction whenever it is well formed.

In both degenerate cases the constrained branch of
Theorem~\ref{thm:keylength} therefore rests \emph{entirely} on
part~(ii), the attack-set licence of A-ATK
(\xref{protocol-compiler/ass:atk}), which no program can verify. Arguably it
is the correct degeneration: in the no-marginal model the licensed class
\emph{is} the scalar class (Remark~\ref{rem:homog}\,(a)), so there is no
wider class for part~(i) to exclude, and the honest statement is that the
attestation is doing all of the work. What an implementation reports of a
passing part~(i) in these cases is Requirement~\ref{req:degenerate-gate-report}.

\emph{The general format.} The scalar class named for the no-marginal model
is the licensed class of the \tref{basic format}. In the general format of
Definition~\ref{def:declaration} the licensed class of
Definition~\ref{def:homog} contains, in \emph{either} mode, the entries
declared \tref{party-local} beside the mode's own form; so in the
no-marginal model the licensed class is the scalar form \emph{together with}
party-local. \emph{And those two are exactly the forms that are
expressible there}: a general right-hand side does not exist in
untrusted-source mode at all (Definition~\ref{def:homog};
Definition~\ref{def:admissible-declaration}\,(vi) in its general-format reading), so it is
refused by well-formedness before part~(i) is reached. \emph{The licensed
class therefore grew by exactly what the expressible class grew by, part~(i)
excludes nothing new, and the degeneration stands}, along the same two
readings: \emph{in the no-marginal model} part~(i) remains a type check that
a declaration passes \emph{by construction} whenever it is well formed, both
expressible forms being licensed --- two forms there, not three; and
\emph{at $d_A=1$} it remains the weaker statement that it rejects nothing a
producer cannot trivially avoid, a general entry being expressible and
unlicensed in trusted-source mode while every $1\times1$ Hermitian
$\Theta_k$ equals $c_k\tau_A$, so that such an entry could always have been
declared homogeneous instead. In both readings the honest statement is still
that the attestation is doing all of the work. In trusted-source mode at
$d_A\ge2$ part~(i) is not degenerate, in either format. \emph{What the third
form does change} is what part~(ii) carries: a tagged entry makes the
attestation attest \emph{more} --- the certificate on the scoped factor's
own channel (\xref{protocol-compiler/rem:party-local-soundness}) --- so the
rules of the gate bind it exactly as they bind the others
(Requirements~\ref{req:constrained-key-gate}
and~\ref{req:degenerate-gate-report}).
\end{remark}
